\documentclass[11pt]{article}
\usepackage[margin=1in]{geometry}
\usepackage{amsmath,amssymb,amsthm,booktabs,microtype}
\usepackage[hidelinks]{hyperref}
\newtheorem{theorem}{Theorem}[section]
\newtheorem{corollary}[theorem]{Corollary}
\newtheorem{remark}[theorem]{Remark}
\newcommand{\dd}{\,\mathrm d}
\title{Normalized Nonlinear Density Comparison on a Gauge Chart\\
\large Separating action oscillation from high-dimensional coverage}
\author{}
\date{20 July 2026}
\begin{document}
\maketitle

\begin{abstract}
We prove an exact event comparison between a reference Gaussian measure and a
nonlinear perturbation on the same retained chart.  The logarithmic loss
separates into remainder oscillation and negative log chart mass.  A centered
Hessian bound $H$ on a radius-$r$ convex ball contributes at most $Hr^2$.
For the side-two four-dimensional block certificate, this local action cost is
only $6.35\times10^{-6}$ at the frozen time.  By contrast, the directly
certified 408-dimensional inner Gaussian ball has log mass $-803.25$, giving a
loss $3.7188$ per plaquette after division over 216 plaquettes.  Thus local
nonlinear geometry is quantitatively adequate, while normalization and chart
coverage remain the unresolved bridge.
\end{abstract}

\section{Normalized comparison theorem}

Let $\nu$ be a probability measure, let $D$ be measurable with $\nu(D)>0$,
and let $\nu_D$ be $\nu$ conditioned on $D$.  Define
\[
 \frac{\dd\mu_D}{\dd\nu_D}(x)=\frac{e^{-R(x)}}{Z_D},
 \qquad Z_D=\int_D e^{-R}\dd\nu_D.
\]

\begin{theorem}[Oscillation and coverage]
If $\operatorname{osc}_D R\leq\kappa$, then for every event $A$,
\begin{equation}
 \mu_D(A)\leq e^\kappa\nu_D(A)
 \leq \exp\!\bigl(\kappa-\log\nu(D)\bigr)\nu(A).
 \label{eq:comparison}
\end{equation}
The analogous lower comparison on events contained in $D$ is
$\mu_D(A)\geq e^{-\kappa}\nu_D(A)$.
\end{theorem}
\begin{proof}
Let $m=\inf_D R$ and $M=\sup_D R$.  Then
$e^{-M}\leq Z_D\leq e^{-m}$, so the normalized density ratio lies between
$e^{-(M-m)}$ and $e^{M-m}$.  This proves the same-chart comparison.  Finally,
$\nu_D(A)=\nu(A\cap D)/\nu(D)\leq\nu(A)/\nu(D)$.
\end{proof}

\begin{corollary}[Centered Hessian certificate]
Suppose $D$ is contained in a convex radius-$r$ ball centered at zero,
$R(0)=0$, $\dd R(0)=0$, and $\|\dd^2R\|\leq H$ on that ball.  Then
\[
 |R(x)|\leq\frac{H}{2}\|x\|^2,
 \qquad \operatorname{osc}_D R\leq Hr^2.
\]
Hence the total log comparison loss in \eqref{eq:comparison} is at most
$Hr^2-\log\nu(D)$.
\end{corollary}
\begin{proof}
Apply Taylor's formula with integral remainder on the segment from zero to
$x$.  The absolute range lies inside $[-Hr^2/2,Hr^2/2]$.
\end{proof}

This formula prevents a common normalization mistake: a pointwise-small
remainder does not compare normalized measures unless the reference mass of
the retained domain is also controlled.

\section{Gaussian chart mass}

For $X\sim N(0,tI_n)$,
\[
 \mathbb P(\|X\|\leq r)=
 P\!\left(\frac n2,\frac{r^2}{2t}\right),
\]
where $P$ is the regularized lower incomplete gamma function.  The executable
certificate evaluates its logarithm by a positive series in the deep lower
tail, avoiding floating-point underflow.

The Stage 08 finite-block input has $136$ scalar physical curvature
directions and three colors, hence $n=408$.  The certified local radius is
$r=0.003$.  Re-evaluating the action certificate at $t=3\times10^{-6}$ gives
$H=0.7051332859$.

\begin{center}
\begin{tabular}{@{}lr@{}}
\toprule
quantity & value\\
\midrule
dimension $n$ & $408$\\
plaquettes & $216$\\
inner radius $r$ & $0.003$\\
action oscillation $Hr^2$ & $6.34620\times10^{-6}$\\
inner-ball log mass & $-803.2526600$\\
total loss per plaquette & $3.71876234$\\
declared target per plaquette & $0.2$\\
\bottomrule
\end{tabular}
\end{center}

The inner Euclidean ball is a certified subset of the local coordinate
domain, so its mass is a valid but potentially very pessimistic coverage lower
bound.  Solving the abstract Gaussian problem gives radius $0.02463$ to meet
the loss target and radius $0.03657$ for log mass $-0.1$.  These numbers are
diagnostics: the nonlinear chart is not certified at those radii.

\section{Frontier consequence}

The action remainder is not the dominant obstruction in this certificate.
One needs a gauge-covariant atlas or enlarged block chart with adequate total
reference mass, controlled overlap multiplicity, transition Jacobians, and no
unjustified global gauge slice.  Only after those data are supplied can
\eqref{eq:comparison} feed a volume-uniform adaptive block estimate.

The theorem is exact on its stated domain but supplies neither full-measure
control nor RG, continuum, or spectral conclusions.  Its claim level is C2.

\end{document}
